\chapter{Sessions: ProPrep remembers every answer}
\label{ch:sessions}

Every time you use ProPrep, it writes down each answer you type, in order, in
a small file. That file is called a \term{session}. It lets you take back a
wrong answer, carry on where you left off after a crash, repeat a whole
preparation without typing, and turn one preparation into a recipe for many
proteins (Chapter~\ref{ch:batch}).

Recording is on from the start. You do not have to do anything to switch it
on.

\section{The session file}
\label{sec:sessions-file}

When ProPrep starts, it prints where it is recording:

\begin{screen}
Session recording to: ./proprep_session_20261121_122825.json
Tip: Use --no-session to disable automatic session recording
Tip: Type 'undo' at any prompt to rewind to an earlier answer
\end{screen}

The name holds the date and the time, so files never overwrite each other. A
short preparation of ubiquitin (download \pdb{1UBQ}, remove the crystal
waters, exit) made a file with 19 answers in it.

\begin{watchout}
The session file is written into the \textbf{project folder}: the folder named
at \prompt{Project name}, not the folder you started ProPrep in. This matters
when you want to resume later (Section~\ref{sec:sessions-resume}).
\end{watchout}

Two answers are never recorded, because they come before recording starts: the
menu-mode question and the project name.

\begin{behindscenes}
The file is plain text in JSON form: a short block of information about the
run (ProPrep's version, the command you typed, the folder, your description if
you gave one) followed by one entry per answer, with the question, your reply
and the menu options you were offered. It is written as you go, so a run that
is interrupted still leaves a usable file.
\end{behindscenes}

You can label a session when you start it, which helps when several files sit
side by side:

\begin{terminal}
proprep --full-menu --session-description "first recording 1UBQ"
\end{terminal}

The text appears in the \prompt{Notes} column when ProPrep later offers you
the file.

\section{Taking back an answer: \texttt{undo}}
\label{sec:sessions-undo}

Type \typed{undo} at any question instead of an answer. ProPrep lists your
recent answers and asks which one to go back to.

\begin{steps}
  \item Type \typed{undo}.
\begin{screen}
Recent answers in this session
 # ┃ module                ┃ question                           ┃ answer
━━━╇━━━━━━━━━━━━━━━━━━━━━━━╇━━━━━━━━━━━━━━━━━━━━━━━━━━━━━━━━━━━━╇━━━━━━━━━━━━━
13 │ PDB Filter            │ Proceed with filtering             │ yes  Yes
14 │ PDB Filter            │ Chain A - Protein retention option │ r  Retain...
15 │ PDB Filter            │ Chain A - Water retention option   │ r  Retain...
16 │ PDB Filter            │ Apply filter selections ...        │ yes  Yes
17 │ PDB Filter Menu       │ Select PDB Filter option           │ 2  Show ...
Rewind to which answer? number to rewind to, [m]ore, or [x] cancel (17): §ans«15»
\end{screen}
  \item Type the number of the answer you want back. \key{m} shows the
        previous page, \key{x} cancels.
  \item ProPrep offers two ways to change it:
\begin{screen}
  15  Chain A - Water retention option  r  Retain entire component
[r]e-answer it live, or [c]hange it and replay what followed? (r): §ans«r»
\end{screen}
  \item Type \key{r}. ProPrep says what it is about to do, replays everything
        up to that question, and then asks it again, live:
\begin{screen}
[Rewinding: question 15 will be asked again; 3 answers from there on are
set aside in a backup of the session log]
...
[Replay complete - now recording new interactions]
\end{screen}
\end{steps}

Answer the question again and carry on as normal. The answers after the
rewind point are kept in a backup file ending in \file{.bak-} plus the time,
in case you want to look at them.

\begin{watchout}
Option \key{c} (\textsf{change it and replay what followed}) does not work
reliably for menu questions in version 1.23.0. We changed a water answer from
\typed{d} to \typed{r} this way; during the replay ProPrep printed
\prompt{[REPLAY: 'r' now means a different option; using 'd' → Discard entire
component]} and kept the old answer. Use \key{r} (re-answer live) instead, or
the Session Editor (Section~\ref{sec:sessions-editor}).
\end{watchout}

\subsection{Where \texttt{undo} does not work}

\begin{itemize}
  \item \textbf{At the project-name question.} Recording has not started, so
        the word is taken literally: a project folder called \file{undo} is
        created.
  \item \textbf{With \texttt{--no-session}.} You get
        \prompt{Invalid choice: undo. Please try again.}
  \item \textbf{During a hands-off replay} (\texttt{--force-replay}), where
        there is nothing to record.
\end{itemize}

Two more things worth knowing. A typing mistake is recorded like any other
answer, so it appears in the list. And rewinding to answer \typed{0} clears
the whole recording: \prompt{[Session truncated to nothing - the first question
will be asked live]}.

\begin{watchout}
In a tool with many settings spread over several menus, check your settings
after an \texttt{undo}. In one of our membrane runs, a value set earlier in a
different sub-menu was back at its default after a rewind
(Chapter~\ref{ch:membrane}).
\end{watchout}

\section{Picking up an earlier session}
\label{sec:sessions-resume}

Start ProPrep \textbf{inside} the project folder that holds the session file.
ProPrep then offers it to you:

\begin{screen}
Previous session files found!

What would you like to do?
  1. Start fresh session (record new interactions)
  2. Resume a previous session from any point
  3. Continue without session recording

Select option  [1/2/3] (1):
\end{screen}

\begin{watchout}
ProPrep only looks in the folder you start it in. Starting it in the folder
\emph{above} the project folder goes straight to the project question, even
though a session file is sitting one level down. Use \texttt{cd} to go into
the project folder first.
\end{watchout}

Choice \key{1} starts a new recording next to the old one. Choice \key{3}
records nothing. Choice \key{2} lists the files, newest first:

\begin{screen}
┏━━━━━┳━━━━━━━━━━━━━━━━━━━━━━━━━━━━━━━━━━━━━━┳━━━━━━━━━━━━━━━━━━┳━━━━━━━━━━━━━━┓
┃ #   ┃ Filename                             ┃ Date/Time        ┃ Interactions ┃
┡━━━━━╇━━━━━━━━━━━━━━━━━━━━━━━━━━━━━━━━━━━━━━╇━━━━━━━━━━━━━━━━━━╇━━━━━━━━━━━━━━┩
│ 1   │ proprep_session_20261121_122825.json │ 2026-11-21 12:28 │ 19           │
└─────┴──────────────────────────────────────┴──────────────────┴──────────────┘

Select session to resume (0 to cancel)  [1/0] (0):
\end{screen}

\begin{watchout}
The default answer here is \key{0}, which means \emph{cancel}. Pressing
\key{Enter} does not resume anything; it starts a fresh session. Type the
number of the file you want.
\end{watchout}

Pick a number and ProPrep offers three ways to use it:

\begin{screen}
Options:
  1. Continue from end (append new interactions)
  2. Open session editor (edit/delete multiple interactions)
  3. Cancel
\end{screen}

Option \key{1} replays every recorded answer and then hands control back to
you. Option \key{2} opens the editor described below. Option \key{3} starts a
fresh session.

\begin{watchout}
\textbf{Do not use \textsf{Continue from end} on a session you ended
normally.} The recording includes your exit, so ProPrep replays that too and
quits before you can type anything:
\begin{screen}
Enter your choice [REPLAY: x]
Are you sure you want to exit? [REPLAY: Yes]
Goodbye!
\end{screen}
Files are re-created on the way through (the structure is downloaded again,
the filter runs again), and then the program is gone. The fix is to delete the
last two answers in the Session Editor first
(Section~\ref{sec:sessions-editor}). \textsf{Continue from end} works as
intended on a session that was interrupted, because those have no recorded
exit.
\end{watchout}

\section{The Session Editor}
\label{sec:sessions-editor}

The editor lets you change or remove earlier answers and then run the
preparation again from the corrected recording. Open it with a flag:

\begin{terminal}
proprep --full-menu --resume-session proprep_session_20261121_122825.json
\end{terminal}

and choose option \key{2}. (The same editor sits behind option \key{2} of the
startup menu.)

\begin{screen}
╭─────────────────────────────── Session Editor ─────────────────────────────╮
│ File: proprep_session_20261121_122825.json | 19 interactions | no changes │
╰───────────────────────────────────────────────────────────────────────────╯
Note: Editing a menu choice that changes the workflow path (e.g., picking
a different tool) may invalidate later interactions. During replay, ProPrep
will fall through to live input when recorded interactions no longer match.

┏━━━━┳━━━━━━━━┳━━━━━━━━━━━━━━━━━━━━┳━━━━━━━━━━━━━━━━━━━━━━━┳━━━━━━━━━━━━━━━━━┓
┃  # ┃ Status ┃ Module             ┃ Question              ┃ Answer          ┃
┡━━━━╇━━━━━━━━╇━━━━━━━━━━━━━━━━━━━━╇━━━━━━━━━━━━━━━━━━━━━━━╇━━━━━━━━━━━━━━━━━┩
│  5 │   ·    │ Structure Loader   │ Enter PDB ID to ...   │ 1UBQ            │
│ 16 │   ·    │ PDB Filter         │ Chain A - Water ...   │ Discard entire  │
│ 18 │   ·    │ PDB Filter Menu    │ Select PDB Filter ... │ Exit            │
│ 19 │   ·    │ Navigation         │ Confirm exit app...   │ Yes             │
└────┴────────┴────────────────────┴───────────────────────┴─────────────────┘
Page 1/1
[e]dit #  [d]elete #|#-#|#-end  [u]ndo #|#-#  [n]ext  [b]ack  [j]ump #
[p]review  save  quit  help
\end{screen}

The answers are shown in words (\prompt{Discard entire component}), not as the
letters you typed. The \prompt{Status} column marks what you have changed:
\texttt{E} for edited, \texttt{D} for deleted, \texttt{·} for untouched.

\subsection{The commands}

\begin{tabularx}{\linewidth}{@{}lL@{}}
\toprule
\textbf{Command} & \textbf{What it does} \\
\midrule
\typed{e 16} & Change answer 16. For a menu question ProPrep lists the
  choices and marks the current one. \\
\typed{d 18}, \typed{d 18-19}, \typed{d 18-end} & Skip these answers when the
  session is replayed. \\
\typed{u 19}, \typed{u 18-19} & Take back a pending change. \\
\typed{n}, \typed{b}, \typed{j 15} & Next page, previous page, jump to the
  page holding answer 15. \\
\typed{p} & List the changes you have made but not yet saved. \\
\typed{save} & Write the corrected file and run it. \\
\typed{quit} & Leave without saving (it asks first if anything is pending).
  In this mode, quitting starts a fresh session. \\
\typed{help} & Show the command list. \\
\bottomrule
\end{tabularx}

\medskip
Changing an answer looks like this:

\begin{screen}
> : §ans«e 16»
Editing #16: Chain A - Water retention option
Available options:
  r: Retain entire component
  s: Select specific residues
  d: Discard entire component <-- current
Current value: d (Discard entire component)
New value (d): §ans«r»
Marked #16 for edit: Discard entire component -> Retain entire component
\end{screen}

\subsection{Removing the recorded exit}

This is the repair for the trap in Section~\ref{sec:sessions-resume}. The last
two answers of a normally-ended session are the exit and the confirmation.
Delete them:

\begin{screen}
> : §ans«d 18-end»
Deleted #18-#19 (2 interactions) (use 'u 18-19' to restore)
\end{screen}

Then type \typed{save}. ProPrep reminds you that a changed answer can send the
run down a different path, asks you to confirm, and writes a new file:

\begin{screen}
Session saved: proprep_session_20261121_122825_edited.json
1 edit(s), 2 deletion(s) applied
...
Hybrid mode started:
  Mode: Replay existing → Record new
\end{screen}

ProPrep then replays the corrected recording and stops at the first question
that is no longer answered, which it now asks you live. Your original file is
not changed; the corrected one gets a new name ending in \file{_edited.json},
and each resume or rewind leaves a \file{.bak-} copy.

\begin{tip}
\prompt{Hybrid mode} is ProPrep's name for ``replay what is recorded, then
record what I type next''. It is the mode you are in after a resume, an edit
or an \texttt{undo}.
\end{tip}

\subsection{When the questions no longer match}

If your change sends the run down a different path, the recorded answers after
it do not fit the new questions. ProPrep does not guess: it stops replaying
and asks you live from that point. The warning at the top of the editor says
so. This is also why editing a tool choice is risky, while editing a value
(a PDB code, a yes/no answer) is usually safe.

\section{Replaying a session, and demonstrating it}
\label{sec:sessions-replay}

To repeat a whole preparation without typing anything, replay the file:

\begin{terminal}
proprep --full-menu --force-replay proprep_session_20261121_122825.json
\end{terminal}

ProPrep prints each question with its recorded answer in brackets and works
through to the end:

\begin{screen}
[Session replay started: proprep_session_20261121_122825.json]
Replaying session from: proprep_session_20261121_122825.json
...
Enter your choice [1a]
Enter PDB ID (e.g., 1ABC, 7XYZ) [1UBQ]
...
Are you sure you want to exit? [Yes]
Goodbye!
[Session replay stopped]
\end{screen}

Our 19-answer session took about 10 seconds and re-created both output files
in the folder we started in. The recording itself is not changed, and no new
session file is written.

For teaching or a demonstration, slow it down so the audience can follow:

\begin{terminal}
proprep --full-menu --force-replay session.json --demo-delay 0.5
\end{terminal}

\begin{screen}
Demo mode: 0.5s delay between actions
\end{screen}

\begin{tip}
Always add \texttt{--full-menu} (or \texttt{--workflow}) to a replay command.
The menu-mode question is not part of a recording, so without the flag ProPrep
stops and asks it before the replay starts.
\end{tip}

The \texttt{--demo-delay} option also works with \texttt{--resume-session}. On
its own it does nothing, and ProPrep says so:
\prompt{Warning: --demo-delay has no effect without --force-replay or
--resume-session}.

\begin{watchout}
A session that was started from a local file with \texttt{--pdbfile} needs
that file again. If it is missing, ProPrep tells you exactly where it looked
and suggests copying the file in or passing \texttt{--pdbfile}, then carries
on with no structure, so every tool reports
\prompt{Needs a loaded structure}.
\end{watchout}

\section{Recovering after a crash or a closed window}
\label{sec:sessions-crash}

Because the file is written as you go, a run that ends badly, by Ctrl-C, a
crash, or a closed browser tab, still leaves everything you had typed. To pick
up:

\begin{terminal}
proprep --full-menu --resume-session proprep_session_20261121_130253.json
\end{terminal}

Choose option \key{1} (\textsf{Continue from end}). ProPrep re-enters each
recorded answer, prints \prompt{[Replay complete - now recording new
interactions]}, and leaves you exactly where you stopped, with recording
continuing in the same file. This is safe here because an interrupted session
has no recorded exit.

We used this for real while writing this manual. A long force-field run
(Chapter~\ref{ch:forcefields}), with dozens of answers in it, was recovered
this way; the replay took seconds.

\begin{watchout}
Resuming replays your \emph{answers}, not ProPrep's memory. It therefore
re-does the work, which may mean downloading and re-running tools. Reopening
the project folder on its own restores nothing: loaded structures and detected
sites are gone until the answers are replayed.
\end{watchout}

\section{Working without a session}

\begin{tabularx}{\linewidth}{@{}lL@{}}
\toprule
\textbf{Option} & \textbf{What it does} \\
\midrule
\texttt{--no-session} & Records nothing and does not offer old sessions. No
  \prompt{Session recording to:} line, and \typed{undo} is refused. \\
\texttt{--session-dir DIR} & Keeps session files in \file{DIR} instead of the
  project folder, and looks there for old ones. \\
\texttt{--session-description TEXT} & A note stored with the file and shown
  in the \prompt{Notes} column. \\
\bottomrule
\end{tabularx}

\medskip
You might turn recording off for a quick look at something, or when you share
a folder and do not want extra files in it. Startup choice \key{3}
(\textsf{Continue without session recording}) does the same for one run.

\begin{watchout}
In version 1.23.0, \texttt{--session-dir} must be an \textbf{absolute} path to
a folder that \textbf{already exists}. A relative path, or a folder that is
not there, stops ProPrep right after the project question with
\prompt{Error: [Errno 2] No such file or directory}, and then confusingly
prints \prompt{Session saved to:} for a file it did not write. The cause is
that the path is worked out after ProPrep has moved into the project folder.
\end{watchout}

\section{If something goes wrong}

\begin{tabularx}{\linewidth}{@{}lL@{}}
\toprule
\textbf{What you see} & \textbf{What to do} \\
\midrule
No \prompt{Previous session files found!} question & You are not in the
  project folder, or you used \texttt{--no-session}. \texttt{cd} into the
  folder that holds the \file{proprep_session_*.json} file. \\
ProPrep replays and exits at once & \textsf{Continue from end} on a session
  that recorded your exit. Delete the last two answers in the editor. \\
\prompt{Invalid choice: undo} & Recording is off for this run. \\
Your edited answer is ignored during replay & Known problem with
  \texttt{undo} option \key{c}. Use \key{r} or the Session Editor. \\
\prompt{Session file not found: nope.json} & Check the name; ProPrep then
  stops. \\
A question is asked live in the middle of a replay & The recording no longer
  matches, usually because the structure differs. Answer it and carry on. \\
\bottomrule
\end{tabularx}

\medskip
One more small oddity: on a normal exit ProPrep prints
\prompt{[Session recording stopped]} but never a \prompt{Session saved to:}
line. The file is there all the same.
